% ============================================================
% JPDC extension --- theory section draft.
% Status labels used throughout:
%   [Definition]  a quantity or measurement boundary is defined;
%   [Derivation]  the statement follows algebraically from the definitions;
%   [Conjecture]  an empirical claim that the W1 design can test only in part.
% No W1 outcome value appears in this file.  All result slots are symbolic.
% ============================================================

\section{End-to-End Energy of Forward-Only Training}
\label{sec:theory}

This section separates the two quantities that are often conflated in
discussions of backpropagation-free learning: the energy of one training pass
and the energy needed to reach a specified accuracy.  The distinction is
elementary, but it is the difference between a per-pass efficiency claim and an
end-to-end efficiency claim.  The W1 family is designed to measure both
quantities under one instrumentation boundary.

\noindent\emph{Status labels.}  A statement labelled \textbf{[Definition]}
fixes a quantity or a measurement boundary; \textbf{[Derivation]} means the
statement follows algebraically from the definitions; \textbf{[Conjecture]}
means the statement is an empirical hypothesis that the W1 design can test only
in part.  Unlabelled sentences are connective exposition and make no
independent empirical claim.  No W1 outcome value appears in this section.

\subsection{Definitions and measurement boundary}

\textbf{[Definition]} Fix a training set $\mathcal{D}$, a model, an
optimisation algorithm, and an evaluation metric $A$ (here, test accuracy).
Let $A^\star$ be a target accuracy fixed before the comparison.  A
\emph{data pass} is one traversal of $\mathcal{D}$ through the training
computation that the algorithm actually performs.  For backpropagation (BP),
one data pass comprises the forward sweep, the backward sweep, and the
parameter update.  For a forward-only algorithm, one data pass comprises every
forward stage, local loss, and local update that the algorithm performs during
that traversal.  This definition deliberately counts the algorithm as
implemented, not an idealised forward FLOP count.

\textbf{[Definition]} Let $e_f$ and $e_b$ be the mean energy consumed per data
pass by the forward-only method and by BP, respectively, measured at the same
device, software stack, batch size, and instrumentation boundary.  Let
$N_f(A^\star)$ and $N_b(A^\star)$ be the number of data passes each method needs
to reach $A^\star$ for the first time.  When the target is unambiguous we write
$N_f$ and $N_b$.

\textbf{[Definition]} The \emph{per-pass energy advantage} of the forward-only
method is
\begin{equation}
  a \;=\; \frac{e_b}{e_f},
  \label{eq:advantage}
\end{equation}
and its \emph{pass multiplier} is
\begin{equation}
  m \;=\; \frac{N_f}{N_b}.
  \label{eq:multiplier}
\end{equation}
Both quantities are dimensionless.  $a>1$ means that one forward-only pass is
cheaper than one BP pass; $m>1$ means that the forward-only method needs more
passes to reach the same target.  Neither statement alone determines the
end-to-end comparison.

\textbf{[Definition]} Ignoring fixed start-up and evaluation costs for the
moment, the end-to-end training energies are
\begin{equation}
  E_f \;=\; N_f\,e_f,
  \qquad
  E_b \;=\; N_b\,e_b .
  \label{eq:endtoend}
\end{equation}
Equations~\eqref{eq:endtoend} are the formal content of the phrase
``end-to-end energy is per-pass energy times the number of passes''.  The
per-pass energies are not assumed equal, and the pass counts are not assumed
equal.

\subsection{The win condition}

\textbf{[Derivation]} From \eqref{eq:advantage}--\eqref{eq:endtoend}, the ratio
of end-to-end energies is
\begin{equation}
  R \;=\; \frac{E_f}{E_b}
      \;=\; \frac{N_f}{N_b}\,\frac{e_f}{e_b}
      \;=\; \frac{m}{a}.
  \label{eq:ratio}
\end{equation}
Therefore, in the no-overhead model,
\begin{equation}
  E_f < E_b
  \quad\Longleftrightarrow\quad
  a > m .
  \label{eq:win}
\end{equation}
A forward-only method with a per-pass advantage $a$ wins end to end exactly
when its pass multiplier is smaller than that advantage.  If $a=m$, the two
methods have equal training energy in this model; if $a<m$, the forward-only
method is cheaper per pass and more expensive end to end.

\textbf{[Derivation]} The result is not a statement about FLOPs.  The
quantities $e_f$ and $e_b$ are measured energies, so \eqref{eq:win}
already includes power draw, achieved clock, memory traffic, and any idle time
inside the pass.  A method can execute fewer arithmetic operations and still
have a smaller $a$ if it utilises the device poorly or moves more data.

\textbf{[Derivation]} Fixed costs change the threshold.  Let
\begin{equation}
  E_f = c_f + N_f e_f,
  \qquad
  E_b = c_b + N_b e_b,
  \label{eq:overhead}
\end{equation}
where $c_f$ and $c_b$ collect start-up, data-loading, evaluation, and other
run-level costs.  Substituting \eqref{eq:advantage}--\eqref{eq:multiplier} and
rearranging gives
\begin{equation}
  a \;>\; \frac{m}
  {1 + \dfrac{c_b-c_f}{N_b e_b}} .
  \label{eq:win-overhead}
\end{equation}
When $c_f=c_b$, \eqref{eq:win-overhead} reduces to \eqref{eq:win}.  A
forward-only method with larger fixed costs must clear a higher per-pass
advantage; a method with smaller fixed costs can tolerate a larger pass
multiplier.  The sign of $c_b-c_f$ therefore matters, and the W1 traces retain
the run boundary and the recorded evaluation timestamps; where a fixed cost
cannot be separated, it is bounded rather than assumed zero.

\textbf{[Definition]} The \emph{matched-accuracy end-to-end energy} of a method
is the cumulative energy consumed from the start of training until the first
pass at which the evaluation metric reaches $A^\star$.  This is the quantity
compared in \eqref{eq:endtoend}.  If a method never reaches $A^\star$, it has no
finite matched-accuracy energy and is reported as a censored run rather than
silently assigned the energy of its last pass.

\textbf{[Conjecture]} Removing the backward sweep is not by itself sufficient
for an end-to-end gain.  The empirically relevant question is whether the
measured pair $(a,m)$ satisfies \eqref{eq:win}; a method that is cheaper per
pass can still lose if its convergence is slow enough.  This conjecture is the
organising hypothesis of the W1 family, not a result of this section.

\subsection{What the W1 design can test}

The cost model is useful only if its quantities are measurable.  Table
\ref{tab:cost-assumptions} lists the assumptions behind
\eqref{eq:win}--\eqref{eq:win-overhead} and the part of the W1 design that can
test each one.  The table is deliberately explicit about what is direct
evidence, what is partial evidence, and what is outside the fixed W1
architecture.

\begin{table}[t]
\centering
\caption{Assumptions of the end-to-end cost model and their W1 tests.  ``Direct''
means the W1 instrumentation records the required quantity; ``partial'' means
the design constrains the assumption but does not vary it; ``not tested'' means
the assumption is a scope limitation.  No outcome value is shown.}
\label{tab:cost-assumptions}
\small
\begin{tabular}{@{}p{0.30\linewidth}p{0.20\linewidth}p{0.42\linewidth}@{}}
\toprule
Assumption & W1 evidence & Comment \\
\midrule
$e_f,e_b$ measured at one boundary
  & direct
  & NVML energy integration with the same start/stop rule for every arm. \\
Same target $A^\star$
  & direct
  & Pre-registered target and evaluation protocol; censored runs reported. \\
$N_f,N_b$ counted as implemented
  & direct
  & Epoch/pass counters; no FLOP-based substitution. \\
Per-pass energy stable over training
  & partial
  & Per-stage sums versus total energy; drift check across seeds. \\
Fixed costs estimated
  & partial
  & Run boundary and evaluation timestamps; fixed costs bounded where not separable. \\
Same data pipeline and batch size
  & direct (design)
  & Frozen configs; no per-arm preprocessing. \\
Same hardware generation
  & partial
  & A100 partition, but node-to-node variation is recorded, not eliminated. \\
Scaling to larger widths/depths
  & not tested
  & W1 fixes the $3\times2000$ MLP; scale is a separate conjecture. \\
\bottomrule
\end{tabular}
\end{table}

\textbf{[Definition]} The \emph{measured} per-pass advantage is
$\hat a=\bar e_b/\bar e_f$, where the bars denote the protocol-level means over
the pre-registered seeds.  The \emph{measured} pass multiplier is
$\hat m=\bar N_f/\bar N_b$ at the same target.  These are estimators of
\eqref{eq:advantage} and \eqref{eq:multiplier}, not definitions of new
quantities.  Their uncertainty must be propagated: a point estimate
$\hat a>\hat m$ is not by itself evidence that $a>m$ in the population of
seeds, because both quantities are measured with run-to-run variation.

\textbf{[Derivation]} If $\hat a$ and $\hat m$ are approximately independent
with relative standard errors $s_a$ and $s_m$, the relative standard error of
$\hat R=\hat m/\hat a$ is approximately
$\sqrt{s_a^2+s_m^2}$.  A decision rule based on \eqref{eq:win} therefore needs
a confidence interval for $R$, not only a comparison of two point estimates.
The W1 analysis reports the interval and the paired-seed structure; the theory
section does not assume that the interval excludes one.

\textbf{[Conjecture]} The pair $(a,m)$ may be architecture-dependent in a way
that the fixed W1 MLP cannot reveal.  In particular, $a$ may improve with width
because activation memory and parameter traffic grow differently in the two
methods, while $m$ may improve with a better local objective or a schedule that
stops earlier.  A scale-dependent crossover, if it exists, would require a
separate width/depth study; W1 can falsify the claim that backward removal is
sufficient at its fixed scale, but it cannot establish a general scaling law.

\subsection{Interpretation rules}

\textbf{[Definition]} We call a result \emph{end-to-end favourable} to the
forward-only method when the matched-accuracy energy ratio satisfies
$R<1$ with a pre-registered uncertainty rule.  We call it \emph{per-pass
favourable} when $a>1$ but make no end-to-end claim from that fact alone.  This
terminology prevents the common slide from ``cheaper per pass'' to ``cheaper
overall''.

\textbf{[Derivation]} Three boundary cases follow directly from
\eqref{eq:ratio}.  First, if $a\le 1$, the forward-only method cannot win end to
end at any $m\ge 1$ in the no-overhead model.  Second, if $a>1$ but
$m\ge a$, the per-pass advantage is exhausted by the extra passes.  Third, if
$m<1$ as well as $a>1$, the forward-only method wins for two independent
reasons; this is a stronger claim than either effect alone and should be
reported separately.

\textbf{[Conjecture]} Early stopping and schedule design act primarily on $m$,
not on $a$.  A fair comparison should therefore report both the stopping rule
and the realised pass count; comparing a forward-only method at its own
convergence with a BP baseline at a fixed epoch budget confounds the two.
The W1 BP-scheduled and budget-matched blocks are designed to separate these
effects, but the conjecture that schedule improvements generalise across
datasets remains to be tested.

\subsection{Summary of the testable model}

The model has one definitional identity, \eqref{eq:endtoend}, one algebraic
win condition, \eqref{eq:win}, and one overhead correction,
\eqref{eq:win-overhead}.  Its empirical content is the pair $(a,m)$ measured at
a fixed target.  The W1 design measures $a$ in the matched-data-pass protocol,
measures $m$ from the pass count to the pre-registered target, retains the
segments needed for $c_f$ and $c_b$, and records node and GPU identity so that
hardware variation is visible rather than absorbed into the treatment effect.
The design does not vary width, depth, or device generation; those remain
conjectures for a later study.
